\uuid{EupT}
\titre{Propriétés de la fonction de répartition}
\auteur{Maxime Nguyen}
\organisation{AMSCC}
\datecreate{2026-11-24}
\niveau{L2}
\module{Probabilité}
\chapitre{Probabilité continue}
\sousChapitre{Espérance et Fonction de répartition}
\difficulte{1}
\theme{fonction de répartition, continuité, dérivabilité}

\contenu{
\begin{question}
Soit $X$ une variable à densité $f$ et $F_X$ sa fonction de répartition. Quelles affirmations sont vraies ?
\begin{answers}
\good{$F_X$ est continue sur $\mathbb{R}$.}
\bad{$F_X$ est strictement croissante sur $\mathbb{R}$.}
\bad{$F_X$ est dérivable en tout point de $\mathbb{R}$.}
\good{Pour tout $t$, $F_X(t) = \mathbb{P}(X < t)$.}
\good{Pour tout $t$, $F_X(t) = \displaystyle\int_{-\infty}^{t} f(x)\,\mathrm{d}x$.}
\end{answers}
\begin{explanations}
B est faux : $F_X$ est seulement croissante. Elle est constante là où $f$ est nulle ; pour la loi uniforme sur $[a;b]$, $F_X = 0$ sur $]-\infty;a]$ et $F_X = 1$ sur $[b;+\infty[$.

C est faux : la fonction de répartition de la loi uniforme sur $[a;b]$ n'est pas dérivable en $a$ ni en $b$. $F_X$ est dérivable sauf éventuellement en un nombre fini de points (pour les densités du cours), et $F_X' = f$ là où $f$ est continue.

D est vrai car $\mathbb{P}(X = t) = 0$.

Contre-sens visés : croire que « à densité » signifie « $F$ partout
dérivable » ou « $F$ strictement croissante » ; refuser D par réflexe du discret.
\end{explanations}
\end{question}
}
